\begin{MyRecipe}{Flammkuchen}{\Calc{3}{\x} Stück}{\SI{30}{\minuteprime} }
	
	\ingredient[\Calc{20}{\x}]{\si{\gram}} {Frischhefe}
	\ingredient[\Calc{1}{\x}]{\si{\Teeloeffel}} {Zucker}
	\ingredient[\Calc{0.4}{\x}]{\si{\kilogram}} {Mehl}
	\ingredient[\Calc{0.2}{\x}]{\si{\liter}} {Wasser}
	\ingredient[\Calc{4}{\x}]{\si{\Essloeffel}} {Öl}
	\ingredient[]{} {Salz}
	
	\Step{Teig}
	Hefe zerbröseln und mit Zucker gut verrühren. Anschließend mit lauwarmen Wasser verquirlen. Mehl mit Salz, Hefe-Wasser und Öl zu einem straffen Teig kneten und \SI{30}{\minuteprime} mit Dampf im Backofen gehen lassen.\par\bigskip
	
	\ingredient[\Calc{1}{\x}]{} {gr. Zwiebel}
	\ingredient[\Calc{0.1}{\x}]{\si{\kilogram}} {Speck}
	
	\Step{Vorbereitung}
	Zwiebeln in feine Halbringe und Speck in feine Streifen schneiden.\par\bigskip

	\ingredient[\Calc{0.5}{\x}]{\si{\kilogram}} {Crème Fraîche}
	\ingredient[\Calc{0.1}{\x}]{\si{\kilogram}} {Gouda}
	\Step{Belegen}
	Teig in gleich große Teile trennen und mit etwas Mehl so dünn wie möglich aus rollen. Ein Teil sollte ausgeollt etwa ein Backofenrost ausfüllen. Backpapiet entsprechend der größe zuschneiden und Teig darauf ausbreiten. Crème Fraîche auf dem Teig verstreichen, dann vorsichtig mit Salz und ordentlich Pfeffer würzen. Dann Zwiebeln und Speck darauf verteilen. Zum Schluss mit etwas Käse bestreuen.\par\bigskip

	\Step{Backen}
	Flammkuchen mit Backpapier vorsichtig auf den Backofenrost ziehen und im vorgeheizen Backofen bei \SI{250}{\celsius} (oder wenn möglich mehr) backen biss er goldbraun ist (circa \SI{10}{\minuteprime}).
	
	
	
\end{MyRecipe}